% Generated draft: requires mathematical and source review.
\documentclass[11pt]{article}
\usepackage[margin=25mm]{geometry}
\usepackage{fontspec}
\setmainfont{Latin Modern Roman}
\usepackage{amsmath,amssymb,mathtools,unicode-math}
\setmathfont{Latin Modern Math}
\usepackage{enumitem,needspace,xurl,hyperref}
\setlength{\parindent}{0pt}
\setlength{\parskip}{5pt}
\setlength{\emergencystretch}{3em}
\newcommand{\N}{\mathbb{N}}
\newcommand{\Z}{\mathbb{Z}}
\newcommand{\Q}{\mathbb{Q}}
\newcommand{\R}{\mathbb{R}}
\newcommand{\C}{\mathbb{C}}
\newcommand{\F}{\mathbb{F}}
\newcommand{\A}{\mathbb{A}}
\newcommand{\PP}{\mathbb{P}}
\newcommand{\E}{\mathbb{E}}
\newcommand{\eps}{\varepsilon}
\newcommand{\dd}{\,\mathrm{d}}
\begin{document}
\section*{UnsolvedMath ID: 20003146}
\textbf{Source problem code:} catalogue entry
\subsection{Definitions and mathematical statement}
The source asks for an application in which an infinite-dimensional Alexandrov space is a limit of manifolds whose dimensions increase. Here ``Alexandrov'' must mean a complete geodesic space satisfying the lower-curvature comparison condition CBB$(0)$ without building finite dimension into the definition. The supplied example is a countable phase configuration space.

Let $(r_j)_{j\geq1}$ be positive real numbers with $\sum_jr_j^2<\infty$. Write each circle of radius $r_j$ as angles $\theta_j\in\R/(2\pi\Z)$ and define
\[
X=\prod_{j\geq1}S^1(r_j),\qquad
 d(\theta,\varphi)^2=\sum_{j\geq1}r_j^2\operatorname{dist}_{\R/(2\pi\Z)}(\theta_j,\varphi_j)^2.
\]
Let $X_N=\prod_{j=1}^NS^1(r_j)$, embedded by fixing all tail angles at zero, and let $p_N:X\to X_N$ be coordinate projection. The aid asserts that $X$ is compact, complete, geodesic, and CBB$(0)$, with infinite dimension since it contains closed tori of every finite dimension. Each $X_N$ is a smooth closed flat $N$-manifold. With
\[
\delta_N=\pi\left(\sum_{j>N}r_j^2\right)^{1/2},
\]
one has $d_H^X(X_N,X)=\delta_N$ exactly. The graph of $p_N$ has distortion at most $\delta_N$, so
\[
d_{\mathrm{GH}}(X_N,X)\leq\delta_N/2\longrightarrow0.
\]
The ambient Hausdorff equality is not asserted to be an exact GH-distance formula.

For the countable $U(1)$-phase/rotor interpretation, let $F:X\to\R$ be $L$-Lipschitz and put $F_N=F|_{X_N}$. Then
\[
|F(x)-F_N(p_Nx)|\leq L\delta_N,
\qquad\left|\min_XF-\min_{X_N}F_N\right|\leq L\delta_N.
\]
Let $\mu$ and $\mu_N$ be normalized product Haar probability measures. The projection coupling gives
\[
\left|\int_XF\,\dd\mu-\int_{X_N}F_N\,\dd\mu_N\right|\leq L\delta_N.
\]
For inverse temperature $\beta>0$, define $Z_\beta=\int_Xe^{-\beta F}\,\dd\mu$ and $Z_{\beta,N}=\int_{X_N}e^{-\beta F_N}\,\dd\mu_N$. Their free energies satisfy
\[
\left|-\beta^{-1}\log Z_\beta+\beta^{-1}\log Z_{\beta,N}\right|\leq L\delta_N.
\]
These are quantitative truncation bounds for fixed Lipschitz observables or energies on this weighted phase space. Normalized Haar measure and a finite uniform Lipschitz constant are essential; arbitrary infinite rotor Hamiltonians or unnormalized product volumes are not covered. The example supplies a mathematical application scheme rather than a claim of a particular physical model or all infinite-dimensional Alexandrov phenomena.
\subsection{Short English statement}
The supplied weighted infinite torus is an infinite-dimensional lower-curvature-comparison space obtained as a compact GH limit of flat finite-dimensional tori. Square-summable radii give explicit tail errors for distances, Lipschitz observables, minima, Haar averages, and finite-temperature free energies. The application requires this weighted metric and normalized probability measures; it does not cover arbitrary infinite rotor energies.
\subsection{Sources}
\begin{itemize}
\item Local registry: permanent catalogue entry, exact source code catalogue entry; \texttt{aim-topology-notes.json}, \texttt{notes[233]}.
\item AIM workshop \emph{Manifolds with nonnegative sectional curvature}, Problem 21, \url{https://aimath.org/WWN/nnsectcurvature/nnsectcurvature.pdf}. The appended ``3. Group Actions and Submersions'' is a following section heading, not part of the mathematical request. No proposer or proposal year is supplied.
\item Supplied research attempt 1, registry updated 14 August 2026: partial result, low novelty confidence, machine-generated and not peer reviewed. It supplies the weighted-circle construction and quantitative phase-space interpretation. Original remarks and literature field are empty; no complete authoritative external references are supplied.
\item Verification reproduces the extraction from \texttt{input.json} but supplies no further definition of infinite-dimensional Alexandrov space or application citation.
\end{itemize}
\end{document}
